\section{Introduction}

Contactless palmprint recognition is a biometric matching problem in which both identification accuracy and low-false-accept verification behavior matter. In practice, a system may be required to retrieve the correct gallery identity while also operating at strict false-acceptance rates under changes in illumination, pose, deformation, region-of-interest quality, and acquisition session. These requirements make the evaluation protocol a central part of the scientific claim. A component that improves performance when development and test identities overlap may not preserve the same ranking when the test palms are held out, the gallery and probe come from different sessions, and the operating point is constrained by a low empirical false-accept rate.

Recent palmprint systems increasingly reuse representation-learning components developed for face recognition and person re-identification, including angular-margin classifiers, supervised contrastive losses, and normalization necks. These components are plausible choices for palmprint embeddings, but their effect is not protocol-independent. In a biometric setting, the relevant question is not only whether a component improves a diagnostic split, but whether the improvement remains under held-out identities or palms, cross-session gallery/probe construction, and conservative verification thresholds. This paper addresses that question through an audited component-level evaluation rather than by proposing a new architecture.

We ask a specific protocol-level question: do conclusions about common recognition heads and losses remain stable when evaluation moves from a seen-identity diagnostic setting to a held-out palm-class cross-session protocol? The question is intentionally narrower than state-of-the-art palmprint recognition. Its importance is methodological: if component rankings change under different split and threshold policies, then claims about recognition-head improvements must be tied to the protocol under which they were measured.

This paper makes four contributions:
\begin{enumerate}
    \item It defines and audits a held-out palm-class Tongji cross-session protocol with bidirectional S1$\rightarrow$S2 and S2$\rightarrow$S1 evaluation, explicit gallery/probe construction, split hashes, and overlap checks.
    \item It evaluates a controlled ResNet18 component matrix covering cross-entropy, supervised contrastive learning, BNNeck, angular-margin supervision (ArcFace, CosFace), and selected combinations under identical checkpoint-selection and matching policies.
    \item It reports conservative low-FAR verification metrics, including TAR@FAR computed only at thresholds whose empirical false-accept rate does not exceed the target.
    \item It analyzes protocol sensitivity: the BNNeck+ArcFace configuration that appears favorable in a seen-identity diagnostic setting does not consistently improve under held-out palm-class cross-session evaluation, and the observed component behavior differs across session direction and operating point.
\end{enumerate}
The paper is therefore a reproducible empirical evaluation, not a new-architecture or state-of-the-art performance claim.
